\documentclass[../../main.tex]{subfiles}
\begin{document}
\section{习题}\label{section:ch22-problems}

\begin{exercise}\label{exercise:ch22-p22-01}
证明\cref{proposition:ch22-subspace-types}（辛子空间、迷向子空间、余迷向子空间和 Lagrang 子空间的性质）。
\end{exercise}

\begin{exercise}\label{exercise:ch22-p22-02}
设 $(V,\omega)$ 是维数为 $2n$ 的辛向量空间。证明：对每个辛、迷向、余迷向或 Lagrang 子空间 $S\subset V$，存在 $V$ 的一组辛基 $(A_i,B_i)$ 具有下列性质：
\begin{enumerate}[(a)]
\item 若 $S$ 是辛的，则 $S=\operatorname{span}(A_1,B_1,\dots,A_k,B_k)$ 对某个 $k$。
\item 若 $S$ 是迷向的，则 $S=\operatorname{span}(A_1,\dots,A_k)$ 对某个 $k$。
\item 若 $S$ 是余迷向的，则 $S=\operatorname{span}(A_1,\dots,A_n,B_1,\dots,B_k)$ 对某个 $k$。
\item 若 $S$ 是 Lagrang 的，则 $S=\operatorname{span}(A_1,\dots,A_n)$。
\end{enumerate}
\end{exercise}

\begin{exercise}\label{exercise:ch22-p22-03}
实辛群是 $\operatorname{Sp}(2n,\mathbb R)\subset\operatorname{GL}(2n,\mathbb R)$，由所有使标准辛张量 $\omega=\sum_i\mathrm{d}x^i\wedge\mathrm{d}y^i$ 不变的 $2n\times2n$ 矩阵组成，即使得对所有 $x,y\in\mathbb{R}^{2n}$ 都有 $\omega(Zx,Zy)=\omega(x,y)$ 的可逆线性映射 $Z:\mathbb{R}^{2n}\to\mathbb{R}^{2n}$ 的集合。
\begin{enumerate}[(a)]
\item 证明矩阵 $Z\in\operatorname{Sp}(2n,\mathbb R)$，当且仅当它把标准基映为辛基。
\item 证明 $Z\in\operatorname{Sp}(2n,\mathbb R)$，当且仅当 $Z^TJZ=J$，其中 $J$ 是 $2n\times2n$ 块对角矩阵
\begin{align*}
J=\begin{pmatrix}j&\cdots&0\\\vdots&\ddots&\vdots\\0&\cdots&j\end{pmatrix},
\end{align*}
主对角线上排列着 $2\times2$ 块 $j=\begin{pmatrix}0&1\\-1&0\end{pmatrix}$，其余位置为零。
\item 证明 $\operatorname{Sp}(2n,\mathbb R)$ 是 $\operatorname{GL}(2n,\mathbb R)$ 的嵌入 Lie 子群，并确定其维数。
\item 确定 $\operatorname{Sp}(2n,\mathbb R)$ 作为 $\mathfrak{gl}(2n,\mathbb R)$ 子代数的 Lie 代数。
\item $\operatorname{Sp}(2n,\mathbb R)$ 是紧的吗？
\end{enumerate}
\end{exercise}

\begin{exercise}\label{exercise:ch22-p22-04}
设 $(M,\omega)$ 是辛流形，$F:N\to M$ 是光滑映射且 $F^*\omega$ 是辛的。证明 $F$ 是光滑浸入。
\end{exercise}

\begin{exercise}\label{exercise:ch22-p22-05}
设 $(M,\omega)$ 是 $2n$ 维紧辛流形。
\begin{enumerate}[(a)]
\item 证明 $\omega^n$（$\omega$ 与自身的 $n$ 重外积）不是恰当的。
\item 证明 $H^{2p}_{\mathrm{dR}}(M)\neq0$，$p=1,\dots,n$。
\item 证明 $S^2$ 是唯一容许辛结构的球面。
\end{enumerate}
\end{exercise}

\begin{exercise}\label{exercise:ch22-p22-06}
证明 $\mathbb{R}^{2n}$（带其标准辛结构）没有任何紧辛子流形。
\end{exercise}

\begin{exercise}\label{exercise:ch22-p22-07}
设 $(M,\omega)$ 与 $(\widehat M,\widehat\omega)$ 是辛流形。在 $M\times\widehat M$ 上定义 $2$-形式 $\Omega=\pi^*\omega-\widehat\pi^*\widehat\omega$，其中 $\pi:M\times\widehat M\to M$、$\widehat\pi:M\times\widehat M\to\widehat M$ 是投影。
\begin{enumerate}[(a)]
\item 证明 $\Omega$ 是辛的。
\item 证明微分同胚 $F:M\to\widehat M$ 是辛微分同胚，当且仅当其图像 $\Gamma(F)=\{(x,y)\in M\times\widehat M:y=F(x)\}$ 是 $(M\times\widehat M,\Omega)$ 的 Lagrang 子流形。
\end{enumerate}
\end{exercise}

\begin{exercise}\label{exercise:ch22-p22-08}
设 $(M,\omega)$ 是辛流形，$S\subset M$ 是余迷向子流形。浸入子流形 $N\subset S$ 如果对每个 $p\in N$ 都有 $T_pN=(T_pS)^\omega$，就称为\emph{特征的}。证明 $S$ 被 $S$ 的连通特征子流形叶状分割，这些叶的维数等于 $S$ 在 $M$ 中的余维数。
\end{exercise}

\begin{exercise}\label{exercise:ch22-p22-09}
把 $\mathbb{R}^{2n}$ 看作带标准辛结构 $\omega=\sum_i\mathrm{d}x^i\wedge\mathrm{d}y^i$ 的辛流形，令 $\Lambda_n\subset G_n(\mathbb{R}^{2n})$ 表示 $\mathbb{R}^{2n}$ 的 Lagrang 子空间所成集合。
\begin{enumerate}[(a)]
\item 证明实辛群 $\operatorname{Sp}(2n,\mathbb R)$ 在 $\Lambda_n$ 上传递作用（见习题 22-3）。
\item 证明 $\Lambda_n$ 有唯一的光滑流形结构使 $\operatorname{Sp}(2n,\mathbb R)$ 的作用光滑，并确定其维数。
\item $\Lambda_n$ 是紧的吗？
\end{enumerate}
\end{exercise}

\begin{exercise}\label{exercise:ch22-p22-10}
证明余切从上的典范辛形式在微分同胚下不变，其意义如下：设 $Q$ 与 $\widehat Q$ 是光滑流形，$F:Q\to\widehat Q$ 是微分同胚。令 $\mathrm{d}F^*:T^*\widehat Q\to T^*Q$ 为习题 11-8 中描述的光滑丛同态。证明当 $T^*Q$ 与 $T^*\widehat Q$ 都赋予各自的典范辛形式时，$\mathrm{d}F^*$ 是辛微分同胚。
\end{exercise}

\begin{exercise}\label{exercise:ch22-p22-11}
设 $Q$ 是光滑流形，$S$ 是 $T^*Q$ 全空间的嵌入 Lagrang 子流形。证明下列命题。
\begin{enumerate}[(a)]
\item 若 $S$ 在某点 $q\in T^*Q$ 处横截于 $T^*Q$ 的纤维，则存在 $q$ 在 $S$ 中的邻域 $V$ 和 $\pi(q)$ 在 $Q$ 中的邻域 $U$，使 $V$ 是定义在 $U$ 上的光滑闭 $1$-形式的像。
\item $S$ 是 $Q$ 上整体定义的光滑闭 $1$-形式的像，当且仅当 $S$ 横截地与每条纤维恰交于一点。（参见定理 6.32 与推论 6.33。）
\end{enumerate}
\end{exercise}

\begin{exercise}\label{exercise:ch22-p22-12}
设 $M$ 是维数至少为 $1$ 的光滑流形。证明：不存在 $M$ 上的 $1$-形式 $\beta$ 使重言形式 $\sigma\in\Omega^1(T^*M)$ 等于拉回 $\pi^*\beta$。
\end{exercise}

\begin{exercise}\label{exercise:ch22-p22-13}
设 $M$ 是光滑流形，$S\subset M$ 是嵌入子流形。定义 $S$ 的\emph{余法丛}为 $T^*M$ 的子集 $N^*S$，定义为
\begin{align*}
N^*S=\{(q,\alpha)\in T^*M:q\in S,\ \alpha|_{T_qS}\equiv0\}.
\end{align*}
证明 $N^*S$ 是 $T^*M|_S$ 的光滑子丛，并且是 $T^*M$（关于 $T^*M$ 上典范辛结构）的嵌入 Lagrang 子流形。
\end{exercise}

\begin{exercise}\label{exercise:ch22-p22-14}
证明下列由 Moser 给出的 Darboux 定理的整体版本 [Mos65]：设 $M$ 是紧光滑流形，$\omega_0$ 是 $M$ 上的辛形式。假设存在光滑含时 $1$-形式 $\alpha:[0,1]\times M\to T^*M$ 使 $\omega_t=\omega_0+\mathrm{d}\alpha_t$ 对每个 $t\in[0,1]$ 都是辛的。证明存在微分同胚 $F:M\to M$ 使 $F^*\omega_1=\omega_0$。
\end{exercise}

\begin{exercise}\label{exercise:ch22-p22-15}
用与\cref{theorem:ch22-darboux}证明中相同的技巧，证明下列 Moser 的定理 [Mos65]：若 $M$ 是可定向紧光滑 $n$ 流形，$n\geqslant1$，$\omega_0,\omega_1$ 是 $M$ 上的光滑定向形式且满足 $\int_M\omega_0=\int_M\omega_1$，则存在微分同胚 $F:M\to M$ 使 $F^*\omega_1=\omega_0$。[提示：对任意定向形式 $\omega$，证明由 $\beta(X)=X\lrcorner\,\omega$（如(16.11)）定义的映射 $\beta:\mathfrak X(M)\to\Omega^{n-1}(M)$ 是丛同构。]
\end{exercise}

\begin{exercise}\label{exercise:ch22-p22-16}
设 $(M,\omega)$ 是辛流形。令 $\mathfrak X(M)\subset\mathfrak X(M)$ 表示 $M$ 上辛向量场所成集合；$\mathfrak H(M)\subset\mathfrak X(M)$ 表示 Hamilton 向量场所成集合。
\begin{enumerate}[(a)]
\item 证明 $\mathfrak X(M)$ 是 $\mathfrak X(M)$ 的 Lie 子代数，$\mathfrak H(M)$ 是 $\mathfrak X(M)$ 的 Lie 子代数。
\item 证明由 $X\mapsto X\lrcorner\,\omega$ 给出的从 $\mathfrak X(M)$ 到 $\Omega^1(M)$ 的映射，诱导出 $\mathfrak X(M)/\mathfrak H(M)$ 与 $H^1_{\mathrm{dR}}(M)$ 之间的向量空间同构。
\end{enumerate}
\end{exercise}

\begin{exercise}\label{exercise:ch22-p22-17}
考虑二体问题，即\cref{example:ch22-n-body}中在 $n=2$ 特殊情形下描述的 Hamilton 系统 $(T^*Q,\omega,E)$。设势能 $V$ 只依赖于质点间的距离。更精确地说，设 $V(q)=v(r(q))$，其中 $v:(0,\infty)\to\mathbb R$ 是某个光滑函数，而
\begin{align*}
r(q)=|q_1-q_2|=\sqrt{(q_1^1-q_2^1)^2+(q_1^2-q_2^2)^2+(q_1^3-q_2^3)^2}.
\end{align*}
\begin{enumerate}[(a)]
\item 设 $u=(u_1,u_2,u_3)$ 是 $\mathbb{R}^3$ 中的单位向量，证明由
\begin{align*}
P(q,p)=u\cdot p_1+u\cdot p_2
=u_1p_1^1+u_2p_1^2+u_3p_1^3+u_1p_2^1+u_2p_2^2+u_3p_2^3
\end{align*}
定义的函数 $P:T^*Q\to\mathbb R$ 是守恒量（称为 $u$ 方向的总线性动量），且对应的无穷小对称生成 $u$ 方向的平移：
\begin{align*}
\theta_t(q,p)=(q_1+tu,q_2+tu,p_1,p_2).
\end{align*}
\item 证明由
\begin{align*}
L(q,p)=q_1^1p_1^2-q_1^2p_1^1+q_2^1p_2^2-q_2^2p_2^1
\end{align*}
定义的函数 $L:T^*Q\to\mathbb R$ 是守恒量（称为关于 $z$ 轴的总角动量），并求出对应无穷小对称的流。说明这与旋转对称性有何关系。
\end{enumerate}
\end{exercise}

\begin{exercise}\label{exercise:ch22-p22-18}
证明\cref{proposition:ch22-conserved-symmetry}（守恒量与无穷小对称的性质）。
\end{exercise}

\begin{exercise}\label{exercise:ch22-p22-19}
本习题给出 Darboux 定理的另一个证明纲要。设 $(M,\omega)$ 是 $2n$ 维辛流形，$p\in M$。
\begin{enumerate}[(a)]
\item 证明开子集 $U\subset M$ 上的光滑坐标 $(x^1,\dots,x^n,y^1,\dots,y^n)$ 是 Darboux 坐标，当且仅当其 Poisson 括号满足
\begin{align}
\{x^i,y^j\}=\delta_{ij},\qquad \{x^i,x^j\}=\{y^i,y^j\}=0.
\end{align}
\item 对 $k$ 作归纳证明下列命题：对每个 $k=0,\dots,n$，存在在 $p$ 处为零、并在 $p$ 的邻域中满足上式的光滑函数 $(x^1,\dots,x^k,y^1,\dots,y^k)$，使 $2k$ 元 $1$-形式组 $(dx^1,\dots,dx^k,dy^1,\dots,dy^k)$ 在 $p$ 处线性无关。当 $k=n$ 时这就证明了定理。[提示：归纳步骤中，假设 $(x^1,\dots,x^k,y^1,\dots,y^k)$ 已找到，找光滑坐标 $(u^1,\dots,u^{2n})$ 使 $\partial/\partial u^i=X_{x^i}$、$\partial/\partial u^{i+k}=X_{y^i}$（$i=1,\dots,k$），并令 $y^{k+1}=u^{2k+1}$。然后找新坐标 $(v^1,\dots,v^{2n})$ 使 $\partial/\partial v^i=X_{x^i}$（$i=1,\dots,k$）、$\partial/\partial v^{i+k}=X_{y^i}$（$i=1,\dots,k+1$），并令 $x^{k+1}=v^{2k+1}$。]
\end{enumerate}
\end{exercise}

\begin{exercise}\label{exercise:ch22-p22-20}
设 $(M,H)$ 是维数为 $2n+1$ 的接触流形。证明：若 $n$ 为奇，则 $M$ 是可定向的；而若 $n$ 为偶，则 $M$ 可定向当且仅当存在 $H$ 的整体接触形式。
\end{exercise}

\begin{exercise}\label{exercise:ch22-p22-21}
证明\cref{theorem:ch22-contact-vector-field}（接触向量场的刻画）。
\end{exercise}

\begin{exercise}\label{exercise:ch22-p22-22}
设 $(M,H)$ 是接触流形，$X$ 是 $M$ 上的光滑向量场。证明：$X$ 是 $H$ 的某个接触形式的 Reeb 场，当且仅当它是一个不在 $H$ 中取值的接触向量场。
\end{exercise}

\begin{exercise}\label{exercise:ch22-p22-23}
证明\cref{theorem:ch22-contact-flowout}（接触流出定理）。
\end{exercise}

\begin{exercise}\label{exercise:ch22-p22-24}
开子集 $U\subset\mathbb R^n$ 上实值函数 $u$ 的经典 eikonal 方程是
\begin{align}
\sum_{i=1}^n\big(\tfrac{\partial u}{\partial x^i}\big)^2=f(x),
\end{align}
其中 $f$ 是给定的光滑实值函数。它在光学理论中起重要作用。（“eikonal”一词源于希腊语中“像”一词，与我们的“icon”一词同根。）在特殊情形 $f(x)\equiv1$ 下，在 $\mathbb R^n$ 的某个开子集上求式(22.35)的一个显式解 $u$，使 $u$ 在单位球面上为零。
\end{exercise}

\begin{exercise}\label{exercise:ch22-p22-25}
设式(22.24)是 Hamilton--Jacobi 方程(22.23)的一个非特征初始条件。对满足(22.25)、(22.26)和(22.27)的任意选取 $\beta:S\to T^*\mathbb R^n$，以及任意 $p\in S$，证明存在 $p$ 的邻域 $U$，使 Cauchy 问题(22.23)--(22.24)在 $U$ 上有唯一解满足对 $x\in S$ 都有 $\mathrm{d}u(x)=\beta(x)$。
\end{exercise}

\begin{exercise}\label{exercise:ch22-p22-26}
证明\cref{theorem:ch22-general-cauchy}（一般一阶 Cauchy 问题的解）。
\end{exercise}
\end{document}
